\documentclass[11pt]{article}

\usepackage{hw_style}

% Set due date for this homework (updated to 2025)
\setduedate{10}{22}

% Setup homework number and header
\hwnumber{4}

\begin{document}

% \title{习题01}
% \author{截止日期: }
% \author{罗凯\\
% 数学物理方法}
% \date{截止日期:}
% \maketitle
 

\begin{hwproblem}{1}
  在指定的点 $z_0$ 的邻域上将下列函数展开为泰勒级数.
  \begin{enumerate}
    \item $\arctan z$的主值, 在 $z_0=0$;(提示:试着利用级数展开然后逐项积分.)
    \item $\ln z$ 在 $z_0=\ii$;
    \item $\ln \left(1+e^z\right)$ 在 $z_0=0$.
  \end{enumerate}
\end{hwproblem}

\vspace{0.5cm}

\begin{hwproblem}{2}
  在指定环域上(或可去除奇点$z_0$处)将下列函数展开为洛朗级数.
  \begin{enumerate}
    \item$1 / z^2(z-1)$ 在 $z_0=1$;
    \item $1 /(z-2)(z-3)$ 在 $|z|>3$;
    \item $\ln \left(1+e^z\right)$ 在 $z_0=0$;
    \item $\frac{1}{z(z+1)}$, 在 $1<|z-i|<\sqrt{2}$.
  \end{enumerate}
\end{hwproblem}

\vspace{0.5cm}

\begin{hwproblem}{3}
 求下列函数在指定点 $z_0$ 的留数:

 \begin{enumerate}
  \item  $\frac{e^{z^2}}{z-1}, z_0=1$;
  \item $\frac{e^{z^2}}{(z-1)^2}, \quad z_0=1$;
  \item $\left(\frac{z}{1-\cos z}\right)^2, z_0=0$;
  \item  $\frac{1+e^z}{z^4}, z_0=0$;
  \item $\frac{1}{z^2 \sin z}, z_0=0$.
 \end{enumerate}

\end{hwproblem}

\vspace{0.5cm}


  
\begin{hwproblem}{4}
分情况讨论 $F(z)=\frac{f^{\prime}(z)}{f(z)}=\frac{d}{d z} \ln f(z)$ 在 $z=a$ 点的性质, 
  若 $a$ 点是 $f(z)$ 的: 
  \\
  (1) $m$ 阶零 点; 
  \\
  (2) $m$ 阶极点.
  \\
  如果 $z=a$ 是 $F(z)$ 的孤立奇点的话, 则求出函数 $F(z)$ 在该点的留数.
\end{hwproblem}

\vspace{0.5cm}


\begin{hwproblem}{5}
  计算下列积分值
  \begin{enumerate}
    \item $\oint_C \frac{d z}{1+z^4}, C$ 为 $|z-1|=1$;
    \item $\oint_{|z|=2} \frac{1}{z^3\left(z^{10}-2\right)} d z$;
    \item $\oint_{|z|=1} \frac{e^z}{z^3} d z$; 
    \item $\int_0^{2 \pi} \frac{d x}{(a+b \cos x)^2}, \quad a>b>0$;
    \item $\int_{-\infty}^{\infty} \frac{d x}{(x^2 + a^2)(x^2+b^2)}, \quad a>0, b>0$;
    \item $\int_{-\infty}^{\infty} \frac{x^2 d x}{\left(x^2+1\right)\left(x^2-2 x \cos \theta+1\right)}, \theta$ 为实数, 且 $\sin \theta \neq 0$.
  \end{enumerate}

\end{hwproblem}

\vspace{0.5cm}

\hwrequirements

\end{document}